\section{Mocninné řady}

\begin{definice}
    Mocninnou řadou rozumíme $\sum_{k=0}^{\infty}c_k\cdot(z-z_0)^k$, kde $c_k \in \mathbb C$ jsou koeficienty řady, $z_0 \in \mathbb C$ je střed řady a $z \in \mathbb C$ je \uv{proměnná}.
\end{definice}

\begin{poznamka}
    Co to je? $\begin{cases}\text{řada s parametrem } z \\ \text{speciální typ funkce} \\ \text{zobecnění polynomu} \end{cases}$
\end{poznamka}

\begin{veta}[Poloměr konvergence] Je dána mocninná řada. Pak $\exists R \in [0, \infty]$ t. ž. $\forall z \in \mathbb C$ platí
\begin{enumerate}[i)]
  \item $|z-z_0|<R \implies$ řada konverguje absolutně,
  \item $|z-z_0|>R \implies$ řada diverguje.
\end{enumerate}
\end{veta}

\begin{proof}
    pomocná množina $\mathscr K \coloneq \left\{\rho \geq 0; \text{ posloupnost } \{|c_k|\cdot \rho^k\}\text{ je omezená}\right\}$

    zřejmě $\mathscr K \neq \emptyset (0 \in \mathscr K)$, položme $R\coloneq \sup \mathscr K$
    \begin{enumerate}[{ad} i)]
        \item nechť $z \in \mathbb C, |z-z_0|<R\implies\exists \rho: \rho >|z-z_0|, \rho \in \mathscr K$

            tj. $\exists C>0$ t. ž. $|c_k|\rho^k\leq C, \forall k$

            lze odhadnout: $|c_k(z-z_0)^k|=|c_k||z-z_0|^k = \underbrace{|c_k|\cdot\rho^k}_{\leq C}\cdot\underbrace{(\frac{|z-z_0|}{\rho})^k}_{\eqcolon \theta^k, \theta <1}\leq C \cdot \rho^k \stackrel{V. 10. 3.}{\implies}$ řada konverguje absolutně
        \item nechť $z \in \mathbb C, |z-z_0|>R\implies\underbrace{|z-z_0|}_{\eqcolon \rho}\notin \mathscr K$

            tj. $\{|c_k||z-z_0|^k\}_{k=0}^\infty$ není omezená

            ?? kdyby řada konvergovala, pak (dle V. 10. 1.) $|c_k(z-z_0)|\to 0, k \to \infty$

            $|c_k(z-z_0)^k|=|c_k||z-z_0|^k$ což není omezené $\implies$ \textsc{spor} \Lightning (V. 7. 2. ?))
    \end{enumerate}
\end{proof}

\begin{poznamka}
    číslo R z předchozí věty se nazývá poloměr konvergence řady. Je určeno jednoznačně.
\end{poznamka}

\begin{proof}
    ?? $\exists R_1 \neq R_2$ splňující vlastnosti i), ii) z V. 11. 1.

    BÚNO $R_1<R_2$ -- zvol $\tilde{z} \in \mathbb C$ t. ž. $R_1<|\tilde{z}-z_0|<R_2$

    pro $\tilde{z}$ řada konverguje i diverguje zároveň -- \textsc{spor} \Lightning
\end{proof}

\begin{veta}
    Je dána mocninná řada. Nechť $c_k \neq 0$, nechť $|\frac{c_{k+1}}{c_k}|\to r, k \to \infty$. Pak $R \coloneq \frac{1}{r}$ (zde rozumíme $\frac{1}{\infty}=0, \frac{1}{0}= \infty$) je poloměr konvergence řady.
\end{veta}

\begin{proof}
    \textsc{trik}: užijeme V. 10. 4. (podíl krit.) na řadu $\sum b_k$, kde $_k = |c_k||z-z_0|^k$ (BÚNO $z\neq z_0$)

    $\frac{b_{k+1}}{b_k}=\frac{|c_{k+1}|}{|c_k|}\cdot \frac{|z-z_0|^{k+1}}{|z-z_0|^k}= \frac{|c_{k+1}|}{|c_k|}\cdot |z-z_0| \to r \cdot |z-z_0|$

      diskuze: \begin{enumerate}[1.]
     \item $\underline{r \in (0, \infty)}: \stackrel{V. 10. 4.}{\implies} \sum b_k$ konverguje, pokud $r\cdot|z-z_0|<1 \implies |z-z_0| < \frac{1}{r}$; $\sum b_k$ diverguje, pokud $r\cdot|z-z_0|>1 \implies |z-z_0| > \frac{1}{r}$

     $\implies$ mocninná řada $\begin{cases}
       \text{konverguje absolutně pro } |z-z_0| <\frac{1}{r},& \\\text{určitě diverguje pro } |z-z_0| >\frac{1}{r},&
       \end{cases}$
     \item $\underline{r=0}: \stackrel{V. 10. 4.}{\implies} \sum b_k$ konverguje $\forall z \in \mathbb C\implies R = \infty$,
     \item $\underline{r= \infty}: b_k \cancel{\to} 0, \forall z \in \mathbb C, z \neq z_0 \implies R = 0.$
   \end{enumerate}
\end{proof}

\begin{veta}
    Je dána mocninná řada. Nechť $c_k \neq 0$, nechť $\sqrt[k]{c_k}\to r, k \to \infty$. Pak $R \coloneq \frac{1}{r}$ (zde rozumíme $\frac{1}{\infty}=0, \frac{1}{0}= \infty$) je poloměr konvergence řady.
\end{veta}

\begin{proof}
  \textsc{trik}: užijeme V. 10. 5. (odmoc. krit.), dále obdobně jak dk. 11. 2.
\end{proof}

\begin{poznamka}
Jedním z hlavních cílů kapitoly je dokázat rovnost
\[
\left\{ \sum_{k=0}^{\infty} c_k (z - z_0)^k \right\}' = \sum_{k=1}^{\infty} k c_k (z - z_0)^{k-1}.
\]
Formálně jde o \uv{derivování řady člen po členu}, neboli o záměnu $\sum$ a $\frac{d}{dz}$.
\end{poznamka}

Řada vpravo je také mocninná řada - můžeme ji přepsat jako
\[
\sum_{k=0}^{\infty} \tilde{c}_k (z - z_0)^k, \quad \tilde{c}_k = (k+1) c_{k+1}.
\]

\begin{lemma}
Řady $(M_1)\sum_{k=0}^{\infty}c_k(z - z_0)^k$ a $(M_2)\sum_{k=1}^{\infty}kc_k(z - z_0)^{k-1}$ mají stejný poloměr konvergence.
\end{lemma}

\begin{dusledek}
Také řady $\sum_{k=2}^{\infty}k(k-1)c_k(z - z_0)^{k-2}$, $\sum_{k=1}^{\infty}\frac{c_k}{k+1}(z - z_0)^{k+1}$ mají stejný poloměr konvergence jako řada $\sum_{k=0}^{\infty}c_k(z - z_0)^k$.
\end{dusledek}

\begin{proof}
  označme $R_1, R_2$ poloměry konvergence řad $(M_1), (M_2)$
  $\stackrel{V. 11. 1.}{\implies} R_1 = \sup \mathscr K_1 \coloneq \left\{\rho \geq 0; \{|\tilde{c}_k|\rho^k\} \text{ je omezená}\right\}, R_2 = \sup \mathscr K_2 \coloneq \left\{\rho \geq 0; \{|\tilde{c}_k|\rho^k\} \text{ je omezená}\right\}$
  \begin{enumerate}[{ad} ]
    \item $\underline{R_2 \leq R_1}: |c_{k+1}|\rho^{k+1} \leq (k+1)|c_{k+1}|\rho^{k+1}$

    $|c_{k+1}|\rho^{k+1} \leq |\tilde{c}_k|\rho^k \cdot \rho,$ tedy pokud $\{\tilde{b}_n\}$ omezená $\implies \{b_k\}$ omezená $\implies \mathscr K_2 \subseteq \mathscr K_1 \implies R_2 \leq R_1$
    \item $\underline{R_2 \geq R_1}:$ ?? $R_2 < R_1$ \dots najděme $\rho_1, \rho_2$ t. ž. $R_2 < \rho_2<\rho_1<R_1$ a $ \rho_2 \notin \mathscr K_2, \rho_1 \in \mathscr K_1$

    $\underbrace{(k+1)|c_{k+1}|\rho_2^k}_{\text{neomezené, neboť } \rho_2 \in \mathscr K_2}=\underbrace{(k+1)\left(\frac{\rho_2}{\rho_1}^k\cdot \frac{1}{\rho_1}\right)}_{\to 0, k \to \infty}\underbrace{|c_{k+1}|\rho_1^{k+1}}_{\leq K\text{, neboť } \rho_1 \in \mathscr K_1}$ -- \textsc{spor} \Lightning
  \end{enumerate}
\end{proof}

\begin{veta}
Nechť řada $\sum_{k=0}^{\infty}c_k(z - z_0)^k$ má poloměr konvergence $R > 0$. Označme
\[
F(z) = \sum_{k=0}^{\infty} c_k (z - z_0)^k, \quad f(z) = \sum_{k=1}^{\infty} k c_k (z - z_0)^{k-1}.
\]
Potom $F'(z) = f(z)$ pro $\forall z \in U(z_0, R)$, kde $U(z_0, R) = \{z \in \C; |z - z_0| < R\}$.
\end{veta}

\begin{proof}
  BÚNO $z_0 = 0, z \in \mathbb C, |z|<R$ pevné

  cíl: $\varphi(h)\to 0, h \to 0$, kde $\varphi(h)= \frac{F(z+h)-F(z)}{h}-f(z)$

  navíc: $0<|h|<\delta; \delta>0$ pevné t. ž. $|z+\delta|<R$

  $F(z+h)=\cancel{c_0}+\cancel{c_1(z+h)}+ \sum_{k=2}^{\infty} c_k (z+h)^k$

  $F(z) = \cancel{c_0} + \cancel{c_1(z)}+\sum_{k=2}^\infty c_k z^k$

  $f(z) = 0 + \cancel{c_1} + \sum_{k=2}^\infty k\cdot c_k \cdot z_{k-1}$

  $\varphi(h) = \sum_{k=2}^\infty c_k \cdot\frac{1}{h}\cdot [\underbrace{(z+h)^k}_{\mathclap{\textstyle =\cancel{z^k}+\cancel{kz^{k-1}h}+\sum_{j=2}^k {k\choose j}z^{k-j}h^j}}-z^k]-kz^{k-1} \stackrel{\text{binom. v.}}{=} \sum_{k=2}^\infty h\cdot c_k\cdot \underbrace{\sum_{j=2}^\infty {k\choose j} z^{h-j}\cdot h^{j-2}}_{\eqcolon S_k}$

  ukážeme, že $\sum_{k=2}\infty |c_k||S_k|<K$ nezávisle na $|h|<\delta$, pak $|\varphi(h)|\leq|h|\cdot K \to 0$

  odhady $S_k: S_k = \sum_{l=0}^{k-2} \underbrace{{k \choose l+2}}_{= \frac{k!}{(l+2)!(k-l-2)!}= \frac{k(k-1)(k-2)!}{\underbrace{(l+2)(l+1)}_{\geq 1}l! (k-l-2)!}}$
\end{proof}

\begin{poznamka}
\begin{itemize}
    \item funkce $F(z)$, $f(z)$ jsou pro $z \in U(z_0, R)$ korektně definovány, neboť dané řady konvergují absolutně (Lemma 11.1.)
    \item tvrzení platí ve smyslu derivace podle komplexní proměnné, tj.
    \[
    (\forall \varepsilon > 0)(\exists \delta > 0)\left[0 < |h| < \delta, h \in \C \Longrightarrow \left| \frac{F(z+h) - F(z)}{h} - f(z) \right| < \varepsilon \right],
    \]
    kde $z \in U(z_0, R)$ je libovolné.
    \item speciální případ je i derivace podle reálné proměnné, tj. $F'(x) = f(x)$, pro každé $x$ z intervalu $(-R, R)$
\end{itemize}
\end{poznamka}

\begin{dusledek}
  Na množině $U(z_0, R)$ platí:
\begin{itemize}
    \item funkce $F(z)$ je nekonečně diferencovatelná a $F''(z) = \sum_{k=2}^{\infty} k(k-1)c_k(z - z_0)^{k-2}$, $F^{(3)}(z) = \sum_{k=3}^{\infty} k(k-1)(k-2)c_k(z - z_0)^{k-3}$ atd.
    \item funkce $\Phi(z) = \sum_{k=0}^{\infty} \frac{c_k}{k+1}(z - z_0)^{k+1}$ je primitivní funkce k $F(z)$, tj. $\Phi'(z) = F(z)$
\end{itemize}
\end{dusledek}

\begin{poznamka}
  Formální definice funkce $e^x$ pomocí mocninných řad (taktéž funkcí $\sin x, \cos x$), položme
  \[
\exp x = \sum_{k=0}^{\infty} \frac{x^k}{k!};
\]
podobně definujeme
\[
\sin x = \sum_{k=0}^{\infty} (-1)^k \frac{x^{2k+1}}{(2k+1)!}, \qquad \cos x = \sum_{k=0}^{\infty} (-1)^k \frac{x^{2k}}{(2k)!},
\]
\end{poznamka}

\begin{lemma}
  Pro $\forall x, y \in \C$ platí
\[
\begin{aligned}
\exp(x + iy) &= \exp(x) \left[ \cos y + i \sin y \right], \\
\sin x &= \frac{1}{2i} \left[ \exp(ix) - \exp(-ix) \right] = \frac{1}{i} \sinh(ix), \\
\cos x &= \frac{1}{2} \left[ \exp(ix) + \exp(-ix) \right] = \cosh(ix).
\end{aligned}
\]
\end{lemma}

\begin{proof}
  TODO
\end{proof}

\begin{veta}
  Nechť řada $F(z) = \sum_{k=0}^{\infty} c_k (z - z_0)^k$ má kladný poloměr konvergence. Potom $F^{(k)}(z_0) = k! c_k$ pro $\forall l = 0, 1, \ldots$.
\end{veta}

\begin{proof}
  TODO
\end{proof}

\begin{poznamka}
  Nechť řada $F(z) = \sum_{k=0}^{\infty} c_k (z - z_0)^k$ má kladný poloměr konvergence. Potom $n$-tý Taylorův polynom funkce $F(z)$ o středu $z_0$ je roven $n$-tému částečnému součtu řady příslušné mocninné řady.
\end{poznamka}

\begin{veta}
  Nechť řady $F(z) = \sum_{k=0}^{\infty} c_k (z - z_0)^k$ a $\tilde{F}(z) = \sum_{k=0}^{\infty} \tilde{c}_k (z - z_0)^k$ mají kladný poloměr konvergence. Nechť $F(z) = \tilde{F}(z)$ na jistém $U(z_0)$. Potom $c_k = \tilde{c}_k$ pro $\forall k = 0, 1, \ldots$.
\end{veta}

\begin{proof}
  TODO
\end{proof}

\begin{poznamka}
  Elementární funkce lze vesměs (lokálně) zapsat mocninnou řadou, kupř.

  $$\ln(1+x) = \sum_{k=0}^\infty (-1)^k \frac{x^{k+1}}{k+1},\, \arctg(x) = \sum_{l=0}^\infty (-1)^l \frac{x^{2l+1}}{2l+1}; x \in (-1, 1).$$
\end{poznamka}

\begin{priklad}
  $C^\infty$ funkce, která není součtem mocninné řady

  $$\phi (x)\coloneq \begin{cases}
    0, &x \leq 0 \\ e^{-\frac{1}{x}}, &x>0
  \end{cases} OBRAZEKOBRAZEK$$

  \begin{itemize}
    \item spojitá v $\mathbb R$ -- slepení dvou spojitých funkcí, v bodě lepení limity sedí
    \item $\phi^\prime (x) = \begin{cases}
      0, &x<0 \\ \frac{1}{x^2}\cdot e^{-\frac{1}{x}}, &x\geq 0\end{cases}$\\
      $\phi^\prime_-(0)=0$\\
      $\phi^\prime_+=0$
    \item obecně $\forall k \geq 0:\phi^{(k)}(0)=0$ -- spojité\\
    ?? $\phi(x)=\sum_{k=0}^\infty c_kx^k, x\in (-R,R), R>0$ poloměr konvergence\\
    $\stackrel{\text{V. 11. 5.}}{\implies} c_k = 0, \forall k \geq 0 \implies \phi(x)\equiv 0$ v $(-R, R),$ což evidentně nemůže být pravda -- \textsc{spor} \Lightning
  \end{itemize}
\end{priklad}
